% 小熊座（综述）
% license CCBYSA3
% type Wiki

本文根据 CC-BY-SA 协议转载翻译自维基百科\href{https://en.wikipedia.org/wiki/Ursa_Minor}{相关文章}。

小熊座（Ursa Minor，拉丁语意为“较小的熊”，与大熊座 Ursa Major 相对），又称小熊星座，是位于遥远北天的一组星座。与大熊座类似，小熊座的尾部也可被视作一只勺子的把手，因此在北美被称为“小北斗”（Little Dipper）：由七颗恒星组成，其中四颗构成勺状的斗部，形态与其对应的大北斗相似。小熊座是公元$^{[2]}$世纪天文学家托勒密列出的$^{[48]}$个星座之一，至今仍是现代$^{[88]}$个星座之一。由于北极星位于其中，小熊座在传统航海导航中，尤其对水手而言，具有重要意义。

北极星（Polaris）是该星座中最亮的恒星，为一颗黄白色超巨星，同时也是夜空中最亮的造父变星，其视星等在$^{[1.97]}$至$^{[2.00]}$之间变化。小熊座$\beta$星（Beta Ursae Minoris），又名 Kochab，是一颗已进入晚期演化阶段的恒星，因膨胀并冷却而成为一颗橙色巨星，视星等为$^{[2.08]}$，仅略暗于北极星。Kochab 与视星等约为$^{[3]}$等的小熊座$\gamma$星（Gamma Ursae Minoris）常被称为“北极星的守护者”或“极星守卫者”$^{[4]}$。在该星座中已有行星被探测到环绕其中四颗恒星运行，包括 Kochab。此外，小熊座还包含一颗孤立的中子星——Calvera，以及迄今发现的第二高温白矮星 H1504+65，其表面温度约为$^{[200{,}000]}$ K。
\subsection{历史与神话}
\begin{figure}[ht]
\centering
\includegraphics[width=6cm]{./figures/9a220c9be61db089.png}
\caption{天龙座环绕着小熊座，如约 1825 年在伦敦出版的一组星座图《乌拉尼娅之镜》（Urania's Mirror）中所描绘的那样。} \label{fig_Ursa_1}
\end{figure}
在巴比伦的星表中，小熊座被称为“天之车”（MUL\,MAR.GÍD.DA.AN.NA，并与女神 Damkina 相关联）。它被列入约公元前 1000 年编纂的《MUL.APIN》星表之中，属于“恩利尔之星”，即北天星区。$^{[6]}$ 据第欧根尼·拉尔修记载并引述卡利马科斯的话，米利都的泰勒斯“测量了腓尼基人航海所依凭的那辆车的星辰”。第欧根尼将其认定为小熊座；由于据说腓尼基人在海上航行时使用该星座，它也被称为 Phoinikē。$^{[7][8]}$ 将北天星座称作“熊”的传统似乎确实源自希腊，尽管荷马只提到过一只“熊”。$^{[9]}$ 因而最初的“熊”应为大熊座；而小熊座作为第二只熊，或“腓尼基之熊”（Ursa Phoenicia，即 Φοινίκη），则是后来才被接纳的。根据斯特拉波（I.1.6, C3）的说法，这是由于泰勒斯的建议，他向希腊人推荐该星座作为航海指引，而此前希腊人一直以大熊座进行航行定位。在古典时代，天极的位置略更接近小熊座的 $\beta$ 星而非 $\alpha$ 星，因此整个星座都被视为指示北方。自中世纪起，使用小熊座 $\alpha$ 星（即“北极星”）作为北极指向变得更加便利，尽管在中世纪北极星仍与天极相距数度。$^{[10][a]}$ 如今，北极星已位于距离北天极 1° 以内，并仍然是当前的极星。其新拉丁名称 \emph{stella polaris} 仅在近代早期才出现。$^{[4]}$  

该星座的另一古名为 Cynosura（希腊语 Κυνοσούρα，“狗之尾”）。这一名称的起源并不明确（若将小熊座视为“狗尾”，则意味着附近应存在一只“狗”的星座，但并无此星座）。$^{[11]}$ 相反，在《星变志》（Catasterismi）的神话传统中，Cynosura 被解释为一位俄瑞阿得斯山林仙女，她曾担任宙斯的乳母，因而受到宙斯的褒奖，被安置于天穹之中。$^{[12]}$ 关于 Cynosura 这一名称还存在多种解释：有一种将其与卡利斯托的神话联系起来，认为其子阿耳卡斯被替换为她的狗，并由宙斯将其置于天上；$^{[11]}$ 也有人认为在更早的观念中，大熊座被视为一头牛，与牧人座构成牧牛组合，而小熊座则是一只狗。$^{[13]}$ 乔治·威廉·考克斯将该名称解释为 Λυκόσουρα 的变体，理解为“狼之尾”，但他本人又将其词源解释为“光的轨迹或行列”（即 λύκος“狼”与 λύκ-“光”之间的关联）。艾伦则指出，可将其与古爱尔兰语中该星座的名称 drag-blod（“火之轨迹”）进行比较。布朗（1899）提出 Cynosura 可能并非希腊语起源，而是借自亚述语 An-nas-sur-ra，意为“高升者”。$^{[14]}$  

另有一种神话叙述两只熊曾在伊达山上将宙斯隐藏起来，使其免遭其父克洛诺斯的杀害；后来宙斯将它们置于天空之中，但由于被神抛向天穹，它们的尾巴被拉得很长。$^{[15]}$  

由于小熊座由七颗恒星组成，拉丁语中表示“北方”的词 \emph{septentrio}（即北极星所指的方向）源自 \emph{septem}（七）与 \emph{triones}（牛），意为“拉犁的七头牛”，而这七颗星也常被想象为这种形象。该名称同样被用来指称大熊座的主星。$^{[16]}$  

在因纽特天文学中，最亮的三颗星——北极星、Kochab 与 Pherkad——被称为 Nuutuittut（“永不移动者”），尽管这一称呼更常以单数形式使用，专指北极星。由于在极北纬地区北极星在天空中的高度过高，它并不适合用于航海定位。$^{[17]}$ 在中国天文学中，小熊座的主要恒星被划分为两个星官：勾陈（Gòuchén，包含 $\alpha$、$\delta$、$\varepsilon$、$\zeta$、$\eta$、$\theta$、$\lambda$ UMi）以及北极（Běijí，包含 $\beta$ 与 $\gamma$ UMi）。$^{[18]}$  
\subsection{特征}
\begin{figure}[ht]
\centering
\includegraphics[width=8cm]{./figures/a72e9261fedb4cea.png}
\caption{约公元 $^{[1009-1010]}$ 年《定星书》中所描绘的小熊座。与西方的表现方式不同，图中的熊被画成尾巴下垂的形态。} \label{fig_Ursa_2}
\end{figure}
小熊座西侧与鹿豹座相邻，西侧亦与天龙座接壤，东侧则与仙王座相邻。其覆盖面积为 $256$ 平方度，在 $88$ 个星座中按面积大小排名第 $56$ 位。小熊座在美国的通俗称呼为“小北斗”，因为其最亮的七颗恒星似乎构成了一个长柄勺（舀勺）的形状。勺柄末端的恒星是北极星 Polaris。北极星也可以通过沿着大熊座“斗碗”末端形成的两颗恒星——大熊座 $\alpha$ 星与 $\beta$ 星（通常称为“指极星”）连成的直线，在夜空中延伸约 $30^\circ$（相当于手臂伸直时三个竖起的拳头的宽度）来找到 $^{[19]}$。构成小北斗“斗碗”的四颗恒星，其视星等分别为二等、三等、四等和五等星，为判断肉眼可见恒星的星等提供了一个简便的参照，对城市居民或用于测试视力尤为有用 $^{[20]}$。

该星座的三字母缩写由国际天文学联合会（IAU）于 $^{[1922]}$ 年正式采用，为“UMi”$^{[21]}$。官方星座边界由比利时天文学家欧仁·德尔波特于 $^{[1930]}$ 年确定，并以一个由 $22$ 条边组成的多边形来界定（见信息框示意）。在赤道坐标系中，这些边界的赤经范围介于 $08^\mathrm{h}\,41.4^\mathrm{m}$ 与 $22^\mathrm{h}\,54.0^\mathrm{m}$ 之间，而赤纬范围则从北天极延伸至南侧的 $65.40^\circ$ $^{[1]}$。由于其位于遥远的北天半球，小熊座整体仅能被北半球的观测者完整看到 $^{[22][b]}$。
\subsection{特征}
\begin{figure}[ht]
\centering
\includegraphics[width=6cm]{./figures/9675489a80e4d9ed.png}
\caption{肉眼所见的小熊座星座（已添加连线与标注）。注意大熊座中构成北斗七星的七颗恒星，并从北斗最外侧的两颗恒星（有时称为“指极星”）连成一条直线指向北极星 Polaris。} \label{fig_Ursa_3}
\end{figure}
\subsubsection{恒星}  
德国制图学家约翰·拜耳使用希腊字母 $\alpha$ 至 $\theta$ 来标注该星座中最显著的恒星，而他的同乡约翰·埃勒特·博德随后又补充使用了 $\iota$ 至 $\phi$。目前仅有 $\lambda$ 与 $\pi$ 仍在使用，这很可能是由于它们靠近北天极的缘故 $^{[16]}$。在该星座的边界之内，共有 $39$ 颗视星等亮于或等于 $6.5$ 等的恒星 $^{[22][c]}$。  

按照约翰·拜耳的排序，主要七颗恒星的传统名称为：  
\begin{enumerate}
\item Polaris  
\item Kochab  
\item Pherkad  
\item Yildun  
\item Epsilon Ursae Minoris 无传统名称。  
\item Zeta Ursae Minoris 无传统名称。  
\item Eta Ursae Minoris 无传统名称。  
\end{enumerate}
作为小熊尾部的标志 $^{[16]}$，北极星 Polaris（即小熊座 $\alpha$ 星）是该星座中最亮的恒星，其视星等在 $1.97$ 与 $2.00$ 之间变化，周期为 $3.97$ 天 $^{[24]}$。它距离地球约 $432$ 光年 $^{[25]}$，是一颗黄白色超巨星，光谱型在 F7Ib 与 F8Ib 之间变化 $^{[24]}$，其质量约为太阳的 $6$ 倍，光度约为太阳的 $2500$倍，半径约为太阳的$45$ 倍。北极星是从地球可见的最亮的造父变星。它是一个三星系统，其中的超巨星主星拥有两颗黄白色主序星伴星，轨道半长轴分别约为 $17$ 和 $2400$ 天文单位（AU），其公转周期分别为 $29.6$ 年与 $42000$ 年 $^{[26]}$。  

传统名称为 Kochab 的小熊座$\beta$星，其视星等为$2.08$，略暗于北极星$^{[27]}$。它距离地球约$131$光年$^{[28][d]}$，是一颗橙色巨星——即核心氢已耗尽、离开主序阶段的演化恒星——光谱型为 K4III$^{[27]}$。Kochab 在约 $4.6$ 天的周期内呈现轻微的光变，其质量通过这些振荡测量估计约为太阳的$1.3$倍$^{[29]}$。Kochab 的光度约为太阳的 $450$倍，直径约为太阳的$42$倍，表面温度约为$4130\,\mathrm{K}$,$^{[30]}$。其年龄估计约为$29.5$亿年，误差约 $\pm 10$ 亿年，并已被宣布拥有一颗行星伴星，其质量约为木星的$6.1$倍，轨道周期为$522$天$^{[31]}$。
\begin{figure}[ht]
\centering
\includegraphics[width=6cm]{./figures/58bc21206ce95031.png}
\caption{小熊座与大熊座相对于北极星 Polaris 的位置关系} \label{fig_Ursa_4}
\end{figure}
传统名称为 Pherkad 的小熊座$\gamma$星，其视星等在 3.04 与 3.09 之间变化，周期约为 3.4 小时$^{[32]}$。它与 Kochab 一同被称为“北极星的守护者”$^{[33]}$。该星是一颗光谱型为 A3II--III 的白色亮巨星 $^{[32]}$，质量约为太阳的 4.8 倍，光度约为太阳的 1050 倍，半径约为太阳的 $15$ 倍 $^{[34]}$，距离地球$487\pm8$光年$^{[28]}$。Pherkad 属于盾牌座 $\delta$ 型变星这一恒星类别$^{[32]}$——这是一类短周期（最长不超过 6 小时）的脉动变星，被用作标准烛光并用于研究星震学 $^{[35]}$。可能同属该类别的还有小熊座 $\zeta$ 星 $^{[36]}$，这是一颗光谱型为 A3V 的白色恒星$^{[37]}$，已开始冷却、膨胀并增亮。它很可能曾是一颗 B3 型主序星，目前呈现轻微变光$^{[36]}$。在视星等 4.95 时，小熊座 $\eta$ 星是小北斗七星中最暗的一颗$^{[38]}$。它是一颗光谱型为 F5V 的黄白色主序星，距离地球 97 光年$^{[39]}$，直径为太阳的 $2$ 倍，质量为太阳的 1.4 倍，光度为太阳的 7.4 倍$^{[38]}$。其附近的小熊座$\zeta$ 星旁还有视星等为 5.00 的小熊座$\theta$星，该星距离地球$860\pm80$ 光年 $^{[40]}$，是一颗光谱型为 K5III 的橙色巨星，已离开主序阶段并发生膨胀与冷却，其直径估计约为太阳的 4.8 倍 $^{[41]}$。

构成小北斗“斗柄”的恒星为小熊座 $\delta$ 星（Yildun）$^{[42]}$与小熊座 $\varepsilon$ 星。$\delta$ 星距离北天极略超过$3.5^\circ$，是一颗光谱型为 A1V 的白色主序星，视星等为 4.35$^{[43]}$，距离地球$172\pm1$光年 $^{[28]}$，其直径约为太阳的 2.8 倍，光度约为太阳的 47 倍$^{[44]}$。小熊座$\varepsilon$星是一个三星系统$^{[45]}$，其合成平均视星等为 4.22 $^{[46]}$。其主星为一颗光谱型为 G5III 的黄色巨星$^{[46]}$，同时也是一颗猎犬座 RS 型变星。该系统为一对光谱双星，伴星距离约 0.36 AU，另有第三颗恒星——一颗光谱型为 K0 的橙色主序星——位于约 8100 AU 处 $^{[45]}$。

靠近北极星的位置有小熊座$\lambda$星，这是一颗光谱型为 M1III 的红巨星。它是一颗半规则变星，视星等在 6.35 与 6.45 之间变化$^{[47]}$。由于该星座位于高北天区，其变星可全年观测：红巨星小熊座 R 星是一颗半规则变星，在 328 天内视星等由 8.5 变至 11.5；而小熊座 S 星则是一颗长周期变星，在 331 天内视星等介于 8.0 与 11 之间$^{[48]}$。位于 Kochab 与 Pherkad 以南、朝向天龙座方向的是小熊座 RR 星$^{[33]}$，这是一颗光谱型为 M5III 的红巨星，同样属于半规则变星，其视星等在 43.3 天内于 4.44 至 4.85 之间变化 $^{[49]}$。小熊座 T 星是另一颗红巨星变星，其状态曾发生显著变化——从视星等在7.8至 15、周期310--315天的长周期（米拉型）变星，转变为半规则变星$^{[50]}$。该恒星被认为经历了一次“壳层氦闪”，即围绕恒星核心的氦壳层达到临界质量并发生点火，这一事件由其在$^{[1979]}$年变光性质的突然改变所标志$^{[51]}$。小熊座 Z 星是一颗微弱的变星，其在$^{[1992]}$年突然变暗$6$个星等，被确认属于一种极为罕见的恒星类别——北冕座 R 型变星$^{[52]}$。

食变星是由于一颗恒星从另一颗前方经过而导致亮度变化的恒星系统，而非源于自身光度变化。小熊座 W 星即为此类系统，其视星等在 1.7 天内于 8.51 与$9.59$ 之间变化$^{[53]}$。该系统的合成光谱型为 A2V，但两颗组成恒星的质量尚不明确。$^{[1973]}$年轨道周期的轻微变化表明该多星系统中可能存在第三个成员——很可能是一颗红矮星——其轨道周期为 $62.2\pm3.9$年$^{[54]}$。小熊座 RU 星是另一例，其视星等在 0.52 天内介于 10 与 10.66 之间 $^{[55]}$。这是一个半分离系统，因为次星充满了其洛希瓣，并正在向主星传递物质$^{[56]}$。

小熊座 RW 星是一个灾变变星系统，其在$^{[1956]}$年爆发为新星，亮度达到$6$等。到$^{[2003]}$年时，其亮度仍比基线值高出 2 个星等，并以每年 0.02 星等的速度逐渐变暗。其距离被计算为$5000\pm800$秒差距（约$16300$光年），这意味着它位于银河系晕之中$^{[57]}$。
“Calvera”这一绰号取自电影《豪勇七蛟龙》中的反派人物，用以指代 ROSAT 全天巡天亮源星表（RASS/BSC）中编号为 1RXS J141256.0+792204 的一处 X 射线源 $^{[58]}$。该天体已被确认是一颗孤立的中子星，是距离地球最近的此类天体之一 $^{[59]}$。小熊座中还包含两颗性质颇为神秘的白矮星。于 $^{[2011\text{-}01\text{-}27]}$ 记录的 H1504+65 是一颗极为暗弱（视星等 $15.9$）的恒星，却拥有目前已知白矮星中最高的表面温度——$200000\,\mathrm{K}$。其大气层由约一半的碳、一半的氧以及 $2\%$ 的氖组成，几乎不含氢和氦，这种成分目前尚无法用现有的恒星演化模型加以解释 $^{[60]}$。WD 1337+705 是一颗温度较低的白矮星，其光谱中含有镁和硅的特征线，这暗示其可能存在伴星或环绕恒星的碎屑盘，但至今尚未发现直接证据$^{[61]}$。WISE 1506+7027 是一颗光谱型为 T6 的棕矮星，距离地球仅约 $11.1^{+2.3}_{-1.3}$ 光年$^{[62]}$。该天体视星等为 14，于 $^{[2011]}$ 年由广域红外巡天探测器（WISE）发现$^{[63]}$。

除 Kochab 之外，还有另外三个恒星系统被发现拥有行星。11 小熊座是一颗光谱型为 K4III 的橙色巨星，其质量约为太阳的 1.8 倍。其年龄约为 15 亿年，自作为 A 型主序星以来已逐渐冷却并发生膨胀。该星距离地球约 390 光年，视星等为 5.04。$^{[2009]}$ 年，人们在其周围发现了一颗质量约为木星 11 倍的行星，其公转周期为 516 天 $^{[64]}$。HD 120084 是另一颗演化恒星，为光谱型 G7III 的黄色巨星，质量约为太阳的 2.4 倍。通过 $^{[2013]}$ 年对该恒星径向速度的精确测量，人们发现其拥有一颗质量约为木星 4.5 倍的行星，该行星具有极高的轨道偏心率（$e=0.66$），是已知偏心率最大的行星轨道之一 $^{[65]}$。HD 150706 是一颗光谱型为 G0V 的类太阳恒星，距离太阳系约 89 光年。最初曾认为其在 $0.6\,\mathrm{AU}$ 处拥有一颗质量与木星相当的行星，但这一说法在 $^{[2007]}$ 年被否定$^{[66]}$。随后于$^{[2012]}$ 年发表的进一步研究表明，该恒星确实存在一颗伴星，其质量约为木星的 2.7 倍，公转周期约为 $16$ 年，轨道半径约为 $6.8\,\mathrm{AU}$,$^{[67]}$。
\subsubsection{深空天体}
\begin{figure}[ht]
\centering
\includegraphics[width=6cm]{./figures/91c3366e922897b8.png}
\caption{NGC\,6217} \label{fig_Ursa_5}
\end{figure}
小熊座相对而言几乎没有深空天体。小熊座矮星系（Ursa Minor Dwarf）是一座矮球状星系，由洛厄尔天文台的阿尔伯特·乔治·威尔逊在$^{[1955]}$年利用帕洛马巡天发现$^{[68]}$。其中心距离地球约 225000 光年$^{[69]}$。在$^{[1999]}$ 年，肯尼思·米格尔与克里斯托弗·伯克利用哈勃空间望远镜证实，该星系只经历过一次恒星形成爆发，发生在约 140 亿年前，并持续了约 20 亿年$^{[70]}$，同时该星系的年龄很可能与银河系本身相当$^{[71]}$。

NGC\,3172（亦称 Polarissima Borealis）是一座极其微弱的星系，视星等为 14.9，恰好是距离北天极最近的 NGC 天体$^{[72]}$。它由约翰·赫歇尔于$^{[1831]}$年发现$^{[73]}$。

NGC\,6217 是一座棒旋星系，距离地球约 6700 万光年$^{[74]}$，使用口径$10\,\mathrm{cm}$（$4\,\mathrm{in}$）或更大的望远镜即可将其定位为一颗约 11 等的天体，位置在小熊座$\zeta$星以东偏北约$2.5^\circ$处$^{[75]}$。该星系被归类为星暴星系，这意味着其恒星形成速率显著高于典型星系 $^{[76]}$。

NGC\,6251 是一座距离地球超过$3.4\times10^8$光年的活动型超巨椭圆射电星系。它拥有一个 Seyfert~2 型活动星系核，并且是 Seyfert 星系中最极端的实例之一。该星系可能与高能伽马射线源 3EG\,J1621+8203 有关联，后者表现出强烈的高能伽马辐射$^{[77]}$。此外，它还以其单侧射电喷流而著称，这是目前已知最明亮的射电喷流之一，于$^{[1977]}$年被发现 $^{[78]}$。

\subsubsection{流星雨}  
天龙座流星雨（Ursids）是一场发生在小熊座方向的显著流星雨，其极大期出现在每年 12 月 18 日至 25 日之间，其母体彗星为 8P/Tuttle$^{[79]}$。
\subsection{另见}  
\begin{itemize}
\item Polaris Flare  
\item 《银河系漫游指南》中的虚构行星 Ursa Minor Beta  
\item 小熊座（中国天文学）
\end{itemize}
\subsection{注释}\\
a.北天极的位置会随着地球自转轴的岁差而发生移动，因此在约 $12000$ 年后，织女星 Vega 将成为新的北极星 $^{[10]}$。\\  
b.尽管对于赤道至南纬 $24^\circ$ 之间的观测者而言，该星座的部分区域在几何上会升出地平线，但靠近地平线数度以内的恒星在实际观测中几乎不可见$^{[22]}$。 \\ 
c.视星等 $6.5$ 的天体属于郊区与乡村过渡夜空条件下肉眼可见的最暗天体之一 $^{[23]}$。\\  
d.更精确地说，其距离为 $130.9\pm0.6$ 光年，该数值由视差测量得到 $^{[28]}$。
\subsection{参考文献}
\begin{enumerate}
\item “小熊座，星座边界”。星座。国际天文学联合会。于2013年6月4日从原文存档。检索日期：2014年5月12日
\item 天文学系（1995）。“小熊座”。威斯康星大学麦迪逊分校。检索日期：2015年6月27日。
\item 柯克帕特里克，J·戴维；马罗科，费德里科等（2024年4月）。“基于全天空20秒差距范围内约3600颗恒星和棕矮星普查的初始质量函数”. 天体物理学报增刊系列. 271 (2): 55. arXiv:2312.03639. 文献标识码:2024ApJS..271...55K. doi:10.3847/1538-4365/ad24e2.
\item "Star Tales – Ursa Minor". Retrieved 2024-07-22.
\item Ridpath, Ian. "Urania's Mirror c.1825 – Ian Ridpath's Antique Star Atlases". Self-published. Retrieved 13 February 2012.
\item 罗杰斯，约翰·H.（1998）。“古代星座的起源：I.美索不达米亚传统”。英国天文协会期刊. 108: 9–28. 文献标识码:1998JBAA..108....9R.
\item 赫尔曼·洪格、大卫·埃德温·平格里，《美索不达米亚的星象科学》 (1999), 第68页.
\item 奥尔布赖特，威廉·F.（1972）。“希腊知识革命中被忽视的因素”。美国哲学学会会刊. 116 (3): 225–42. JSTOR 986117.
\item 里德帕斯，伊恩。“小熊座”。星图。自费出版。检索日期：2015年3月7日布洛姆伯格，彼得·E.（2007）。“星座‘大熊座’为何得名？”（PDF）。载于帕斯托尔，埃米莉亚（编）。考古天文学在考古学与民族志中的应用：2004年于匈牙利凯奇凯梅特举行的欧洲文化天文学学会年度会议论文集。英国牛津：Archaeopress。第129–32页。ISBN978-1-4073-0081-8。归档于原文2016-03-282015-08-01
\item 肯尼思·R·朗（2013）。基础天体物理学。施普林格科学与商业媒体。第10–15. ISBN 978-3-642-35963-7.
\item 艾伦，理查德·欣克利（1899）。恒星名称：其传说与含义. 447页。“这个词的词源尚不确定，因为这一星群并不完全符合其名称，除非将狗本身也纳入其中；不过，仍有一些人联想到卡利斯托与其犬而非阿尔卡斯的另一种传说，认为这或许正是答案所在。”其他人则将这一名称源自马拉松以东的阿提卡海角，因为水手们从海上驶近该海角时，会看到这些星辰在其上方及更远处闪耀；但若二者之间真有何种关联，倒推其来源似乎更为合理；而博尔努夫则坚称，它与希腊语中“狗”一词毫无关联。
\item 康多斯，T.，《卡塔斯特里莫伊》（第一部分），1967年。此外，还被塞尔维乌斯提及于维吉尔的《农事诗》1.246，约公元400年；关于其真实性存疑的记载见希吉努斯的《天文学》2.2。
\item 265f. 罗伯特·布朗关于希腊、腓尼基和巴比伦原始星座起源的研究（1899年）, “M. 西罗诺斯（古代钱币类型集第116页）认为，在某些克里特硬币图案中，大熊座被描绘为一头母牛，因此Boôtês意为“牧人”，而小熊座则被视作一只狗（“Chienne” 参见Kynosoura、Kynoupês），即宙斯的哺乳者。”据称，拉丁传统中将小熊座命名为Catuli‘幼崽’或Canes Laconicae‘斯巴达犬’，这一说法记载于约翰·海因里希·阿尔施泰德（1649年，第408页）之中，很可能是一种近代早期的创新。
\item 然而，就在不久前，布朗［罗伯特·布朗，《希腊、腓尼基与巴比伦原始星座起源研究》］提出，这一词并非源自希腊语，而是源于幼发拉底河流域；为佐证这一观点，他提及了该流域中一个星座的名称，赛斯将其转写为An‑ta-sur‑ra，意为“上层天球”。布朗则将其读作An‑nas-sur‑ra，即“高高升起”，这一译法与小熊座极为贴切。他还将其与Κ‑υν‑όσ‑ου‑ρα相比较——若略去首字母辅音，则变为Unosoura。（艾伦，理查德·欣克利，《星名：其传说与含义》，纽约，多佛出版社，1963年，第448页）布朗指出，阿拉图斯恰如其分地将“Cynosura”形容为“高高运行”（“夜深人静时，犬尾座的头部高高升起”，κεφαλὴ Κυνοσουρίδος ἀκρόθι νυκτὸς ὕψι μάλα τροχάει，第308行及以下）。
\item 罗杰斯，约翰·H.（1998）。“古代星座的起源：II. 地中海传统”。英国天文协会期刊. 108: 79–89. 文献标识码:1998JBAA..108...79R.
\item 瓦格曼，莫顿（2003）。失落的星辰：约翰内斯·拜耳、尼古拉斯·路易·德·拉凯耶、约翰·弗莱姆斯特德及其他多位天文学家星表中的失踪、遗失与疑难恒星。弗吉尼亚州布莱克斯堡：麦当劳与伍德沃德出版公司。第312页，第518页。ISBN 978-0-939923-78-6.
\item 麦克唐纳，约翰（1998）。《北极天空：因纽特人的天文学、星象传说与神话》。安大略省多伦多：皇家安大略博物馆/努纳武特研究所。第61. ISBN 978-0-88854-427-8.
\item “小熊座——中国关联”. 星图故事。检索日期：2023-01-13.
\item 奥米拉，斯蒂芬·詹姆斯（1998）。梅西耶天体。深空伴侣。英国剑桥：剑桥大学出版社。第10. ISBN 978-0-521-55332-2.
\item 奥尔科特，威廉·泰勒 (2012) [1911]. 各时代星象传说：北半球星座的神话、传说与事实集。纽约，纽约：信使公司。第377页。ISBN 978-0-486-14080-3.
\item 拉塞尔，亨利·诺里斯 (1922). “星座的新国际符号”. 大众天文学. 30: 469. 文献标识码:1922PA.....30..469R.
\item 里德帕斯，伊恩。“星座：蜥蜴座—狐狸座”。星图故事。自费出版。检索日期：2014年6月21日
\item 博特尔，约翰·E.（2001年2月）。“博特尔暗空等级”。天空与望远镜。于2014年3月31日从原文存档。检索于2014年11月29日
\item 奥特罗，塞巴斯蒂安·阿尔贝托（2007年12月4日）。“小熊座α星”。国际变星索引。美国变星观测者协会。检索日期：2014年5月16日
\item “小熊座α星——经典造父变星（δ 仙王型）”。SIMBAD天文学数据库。斯特拉斯堡天文数据中心。检索日期：2014年8月19日
\item 卡勒，詹姆斯·B.“北极星”。恒星。伊利诺伊大学。检索日期2014年8月19日
\item “小熊座β——变星”。SIMBAD天文学数据库。斯特拉斯堡天文数据中心。检索日期：2014年5月18日
\item 范·勒文，F.（2007）。“新依巴谷归算的验证”。天文学与天体物理学. 474 (2): 653–64. arXiv:0708.1752. 文献编码:2007A&A...474..653V. doi:10.1051/0004-6361:20078357. S2CID 18759600.
\item 塔兰特，新泽西州；查普林，W.J.；埃尔沃思，Y.；斯普雷克利，S.A.；史蒂文斯，I.R.（2008年6月）。“大熊座β星的振荡：利用SMEI的观测结果”。天文学与天体物理学.483(第3期)：L43－L46arXiv：0804.3253文献编码2008A&A...483L..43Tdoi10.1051/0004-6361:200809738S2CID53546805
\item 卡勒，詹姆斯·B.“科恰布”。恒星。伊利诺伊大学。检索日期：2014年8月19日
\item 李B.-C.；韩I.；朴M.-G.；姆克里奇安D.E.；哈策斯A.P.；金K.-M.（2014）。“位于巨蟹座β、狮子座μ和小熊座β的类木行星伴星”。天文学与天体物理学. 566: 7. arXiv:1405.2127. Bibcode:2014A&A...566A..67L. doi:10.1051/0004-6361/201322608. S2CID 118631934. A67。
\item 沃森，克里斯托弗（2010年1月4日）。“小熊座γ”。国际变星索引。美国变星观测者协会。检索日期：2014年5月18日
\item 阿诺德，H. J. P.；多赫蒂，保罗；摩尔，帕特里克（1999）。恒星摄影图集。佛罗里达州博卡拉顿：CRC出版社。第148页。ISBN 978-0-7503-0654-6.
\item 卡勒，詹姆斯·B.（2013年12月20日）。“珀尔卡德”。恒星。伊利诺伊大学。检索日期：2014年5月18日
\item 坦普尔顿，马修（2010年7月16日）。“盾牌座δ星与盾牌座δ变星”。当季变星。美国变星观测者协会。检索日期：2014年8月19日
\item 卡勒，詹姆斯·B.“阿尔法·法尔卡代因”.恒星. 伊利诺伊大学。检索日期：2014年6月21日
\item “小熊座ζ星——变星”。SIMBAD天文学数据库。斯特拉斯堡天文数据中心。检索日期：2014年6月21日
\item 卡勒，詹姆斯·B.“安瓦尔·阿尔·法尔卡代因”.星辰. 伊利诺伊大学. 检索日期2014年6月21日
\item “小熊座η”.SIMBAD天文学数据库。斯特拉斯堡天文数据中心。检索日期2014年7月30日
\item “小熊座θ——变星”。SIMBAD天文学数据库。斯特拉斯堡天文数据中心。检索日期：2014年7月30日
\item 帕西内蒂·弗拉卡西尼，L. E.；帕斯托里，L.；科维诺，S.；波齐，A.（2001年2月）。“恒星视直径与绝对半径目录（CADARS）——第三版——注释与统计”。天文学与天体物理学. 367 (2): 521–24. arXiv:astro-ph/0012289. 文献编码:2001A&A...367..521P. DOI:10.1051/0004-6361:20000451. S2CID 425754.
\item “命名恒星”. IAU.org。于2025年3月10日从原文存档。检索日期：2018年8月8日。
\item “小熊座δ”.SIMBAD天文学数据库。斯特拉斯堡天文数据中心。检索日期2014年6月21日
\item 卡勒，詹姆斯·B.“耶尔顿”.星辰. 伊利诺伊大学. 检索日期2014年7月30日
\item 卡勒，詹姆斯·B.“小熊座ε星”.恒星. 伊利诺伊大学。检索日期2014年6月21日
\item “小熊座ε——RS CVn型变星”。SIMBAD天文学数据库。斯特拉斯堡天文数据中心。检索日期：2014年6月21日
\item 沃森，克里斯托弗（2010年1月4日）。“小熊座λ星”。国际变星索引。美国变星观测者协会。检索日期：2014年6月21日
\item 莱维，大卫·H.（1998）。观测变星：初学者指南。英国剑桥：剑桥大学出版社。第133页。ISBN 978-0-521-62755-9.
\item 奥特罗，塞巴斯蒂安·阿尔贝托（2009年11月16日）。“小熊座RR星”。国际变星索引。美国变星观测者协会。检索日期：2014年5月18日
\item 乌滕塔勒，S.；范斯蒂普豪特，K.；沃特，K.；范温克尔，H.；范埃克，S.；约里森，A.；克尔施鲍姆，F.；拉斯金，G.；普林斯，S.；佩塞米耶，W.；瓦埃尔肯斯，C.；弗雷马，Y.；亨斯贝格，H.；杜莫蒂耶，L.；莱曼，H.（2011）。“脉动周期变化的米拉变星的演化状态”。天文学与天体物理学. 531: A88.arXiv:1105.2198. 文献标识码:2011A&A...531A..88U. DOI:10.1051/0004-6361/201116463. S2CID 56226953.
\item 马泰伊，珍妮特·A.；福斯特，格兰特（1995）。“大熊座T星的戏剧性周期性亮度下降”。美国变星观测者协会期刊. 23 (2): 106–16. 天文学目录编码:1995JAVSO..23..106M.
\item 本森，普里西拉·J.；克莱顿，杰弗里·C.；加尔纳维奇，彼得；斯科迪，宝拉（1994）。“Ursa Minoris——一颗新的北冕座R型变星”. 天文学杂志. 108 (#1): 247–50. 文献标识码:1994AJ....108..247B. DOI:10.1086/117063.
\item 沃森，克里斯托弗（2010年1月4日）。“小熊座W星”。国际变星索引。美国变星观测者协会。检索日期：2015年7月18日
\item 克雷纳，J. M.；普里布拉，T.；特雷姆科，J.；斯塔霍夫斯基，G. S.；扎克热夫斯基，B. （2008）。“三颗近双星系统的周期分析：TW 安德、TT 赫拉和W 小熊座”. 皇家天文学会月报. 383 (#4): 1506–12. 文献编码:2008MNRAS.383.1506K. DOI:10.1111/j.1365-2966.2007.12652.x.
\item 沃森，克里斯托弗（2010年1月4日）。“小熊座RU星”。国际变星索引。美国变星观测者协会。检索日期：2015年7月18日
\item 马尼马尼斯，V. N.；尼亚尔霍斯，P. G.（2001）。“近接触双星系统小熊座RU的光度研究”. 天文学与天体物理学. 369 (3): 960–64. 文献标识码:2001A&A...369..960M. doi:10.1051/0004-6361:20010178.
\item 比安基尼，A.；塔佩特，C.；坎特纳，R.；坦布里尼，F.；奥斯本，H.；坎特雷尔，K.（2003）。“小熊座RW星（1956）：一个演化的新星后系统”. 太平洋天文学会出版物. 115 (#809): 811–18. 文献编码:2003PASP..115..811B. DOI:10.1086/376434.
\item “在地球附近发现罕见的死亡恒星”.BBC新闻：科学/自然. BBC. 2007年8月20日.已存档自原始网页，日期为2014年7月13日。检索日期2007年8月21日
\item 拉特利奇，罗伯特；福克斯，德里克；谢夫丘克，安德鲁（2008）。“在高银纬处发现一颗孤立的致密天体”。天体物理学报. 672 (#2): 1137–43. arXiv:0705.1011. 文献编码:2008ApJ...672.1137R. DOI:10.1086/522667. S2CID 7915388.
\item 维尔纳，K.；劳赫，T.（2011）。“利用哈勃太空望远镜宇宙起源光谱仪对炽热裸星核H1504+65的紫外光谱研究”。天体物理学与空间科学. 335 (1): 121–24. 文献标识码:2011Ap&SS.335..121W. DOI:10.1007/s10509-011-0617-x. S2CID 116910726.
\item 迪金森，N. J.；巴尔斯托，M. A.；韦尔什，B. Y.；伯利，M.；法里希，J.；雷德菲尔德，S.；昂格劳布，K.（2012）。“炽热白矮星周围特征的起源”. 皇家天文学会月报. 423 (2): 1397–1410. arXiv:1203.5226. 文献标识码:2012MNRAS.423.1397D. DOI:10.1111/j.1365-2966.2012.20964.x. S2CID 119212643.
\item 马什，肯尼斯·A.；赖特，爱德华·L.；柯克帕特里克，J·戴维；杰利诺，克里斯托弗·R.；库欣，迈克尔·C.；格里菲斯，罗杰·L.；斯克鲁茨基，迈克尔·F.；艾森哈特，彼得·R.（2013）。“光谱类型为Y和晚期T的超冷棕矮星的视差与自行”。天体物理学报. 762 (2): 119. arXiv:1211.6977. Bibcode:2013ApJ...762..119M. doi:10.1088/0004-637X/762/2/119. S2CID 42923100.
\item 柯克帕特里克，J. 戴维；库欣，迈克尔 C.；盖利诺，克里斯托弗 R.；格里菲斯，罗杰 L.；斯克鲁茨基，迈克尔 F.；马什，肯尼斯 A.；赖特，爱德华 L.；梅因泽，艾米 K.；艾森哈特，彼得 R.；麦克莱恩，伊恩 S.；汤普森，玛姬 A.；鲍尔，詹姆斯 M.；本福德，多米尼克 J.布里奇，凯莉·R.；莱克，肖恩·E.；佩蒂，萨拉·M.；斯坦福，斯宾塞·亚当；蔡朝伟；贝利，凡妮莎；贝希曼，查尔斯·A.；布鲁姆，乔舒亚·S.；博昌斯基，约翰·J.；伯格瑟，亚当·J.；卡帕克，彼得·L.；克鲁兹，凯勒·L.；欣茨，菲利普·M.；卡尔塔尔特佩，杰伊汉·S.；诺克斯，罗素·P.；马诺哈尔，斯瓦尔尼玛；马斯特斯，丹尼尔；莫拉莱斯-卡尔德龙，玛丽亚；普拉托，丽莎·A.；罗迪加斯，蒂莫西·J.；萨尔瓦托，玛拉；舒尔，史蒂文·D.；斯科维尔，尼古拉斯·Z.；辛姆科，罗伯特·A.；斯塔佩尔费尔特，卡尔·R.；斯特恩，丹尼尔；斯托克，内森·D.；瓦卡，威廉·D. （2011）。“宽视场红外巡天探测器（WISE）发现的首批一百颗棕矮星”。天体物理学杂志增刊。197（2）：19。arXiv：1108.4677v1文献标识码2011ApJS..197...19KDOI10.1088/0067-0049/197/2/19S2CID16850733
\item 多林格，M. P.；哈策斯，A.P.；帕斯奎尼，L.；根特尔，E. W.；哈特曼，M.（2009）。“围绕K型巨星小熊座11和HD 32518的行星伴星”。天文学与天体物理学. 505 (3): 1311–17. arXiv:0908.1753. 文献编码:2009A&A...505.1311D. doi:10.1051/0004-6361/200911702. S2CID 9686080.
\item 佐藤文英；大宫正志；原川宏树；刘玉娟；泉浦秀幸；神部荣二；武田洋一；吉田道敏；伊藤洋一；安藤博康；国分荣一郎；伊田茂（2013）。“三颗演化中的中等质量恒星的行星伴星：HD 2952、HD 120084 和蛇夫座ω”. 日本天文学会出版物. 65 (4): 1–15. arXiv:1304.4328. Bibcode:2013PASJ...65...85S. doi:10.1093/pasj/65.4.85. S2CID 119248666.
\item 赖特，J.T.；马西，G.W.；费舍尔，D.A.；巴特勒，R.P.；沃格特，S.S.；廷尼，C.G.；琼斯，H.R.A.；卡特，B.D.；约翰逊，J.A.；麦卡锡，C.；艾普斯，K.（2007）。“四颗新系外行星及系外行星宿主恒星可能存在更多次级伴星的迹象。”天体物理学报. 657 (1): 533–45. arXiv:astro-ph/0611658. Bibcode:2007ApJ...657..533W. doi:10.1086/510553. S2CID 35682784.
\item 博瓦斯，伊莎贝尔；佩佩，弗朗切斯科；佩里耶，克里斯蒂安；克洛兹，迪迪埃；邦菲尔斯，哈维埃；布希，弗朗索瓦；桑托斯，努诺·C.；阿诺德，吕克；博济，让-吕克；迪亚斯，罗德里戈·F.；德尔福斯，哈维埃；埃根贝格尔，安妮；埃伦赖希，大卫；福尔韦耶，蒂埃里；埃布拉尔，纪尧姆；拉格朗日，安妮-玛丽；洛维斯，克里斯托夫；马约尔，米歇尔；穆图，克莱尔；纳埃夫，多米尼克；桑特讷，亚历山大；塞格朗桑，达米安；西万，让-皮埃尔；于德里，斯蒂芬（2012），“SOPHIE北天系外行星搜寻Ⅴ：对ELODIE候选行星的后续观测——类木星行星围绕类太阳恒星运行”, 天文学与天体物理学, 545: A55,arXiv:1205.5835, 文献标识码:2012A&A...545A..55B, DOI:10.1051/0004-6361/201118419, S2CID 119109836
\item 伯格，西德尼（2000）。本星系群的星系。英国剑桥：剑桥大学出版社。第257页。ISBN 978-1-139-42965-8.
\item 格雷贝尔，伊娃·K.；加拉格尔，约翰·S.三世；哈贝克，丹尼尔（2003）。“矮椭圆星系的祖先”。天文学杂志. 125 (4): 1926–39. arXiv:astro-ph/0301025. Bibcode:2003AJ....125.1926G. doi:10.1086/368363. S2CID 18496644.
\item 范登伯格，西德尼（2000年4月）：“本星系群的最新信息”。太平洋天文学会出版物. 112 (#770): 529–36. arXiv:astro-ph/0001040. 文献编码:2000PASP..112..529V. DOI:10.1086/316548. S2CID 1805423.
\item 米格尔，肯尼思·J.；伯克，克里斯托弗·J.（1999）。“大熊座矮椭圆星系的WFPC2观测”。天文学杂志. 118 (366): 366–380. arXiv:astro-ph/9903065. 文献标识码:1999AJ....118..366M. DOI:10.1086/300923. S2CID 119085245.
\item “NGC 3172”. sim-id. 检索日期2020-05-29.
\item “新总目录天体：NGC 3150 - 3199”. cseligman.com。检索日期：2020-05-30.
\item 古谢夫，A. S.；皮柳金，L. S.；萨希波夫，F.；多东诺夫，S. N.；叶日科娃，O. V.；赫拉姆佐娃，M. S.；加尔松胡赫德，F.（2012）。“六个旋涡星系中H II区的氧和氮丰度”. 皇家天文学会月报. 424 (#3): 1930–40. arXiv:1205.3910. 文献标识码:2012MNRAS.424.1930G. DOI:10.1111/j.1365-2966.2012.21322.x. S2CID 118437910.
\item 奥米拉，斯蒂芬·詹姆斯（2007）。史蒂夫·奥米拉的赫歇尔400观测指南。英国剑桥：剑桥大学出版社。第227页。ISBN 978-0-521-85893-9.
\item 卡尔泽蒂，丹妮拉（1997）。“星暴星系中的红化与恒星形成”。天文期刊. 113: 162–84. arXiv:astro-ph/9610184. Bibcode:1997AJ....113..162C. doi:10.1086/118242. S2CID 16526015.
\item “NGC 6251 – 赛弗特2型星系”.SIMBAD天文学数据库. 斯特拉斯堡天文数据中心. 检索日期：2015年7月21日
\item 珀利，R. A.；布里德尔，A. H.；威利斯，A. G.（1984）。“对NGC 6251中射电喷流的高分辨率VLA观测”。天体物理学杂志增刊系列. 54: 291–334. 文献标识码:1984ApJS...54..291P. doi:10.1086/190931.
\item 詹尼斯肯斯，彼得（2012年9月）。“绘制流星体轨道：新流星雨被发现”。天空与望远镜: 24.
\end{enumerate}
\subsection{外部链接}
\begin{itemize}
\item 深空摄影星座指南：小熊座
\item 可点击的小熊座
\item 瓦尔堡学院图像数据库（约160幅中世纪及早期现代的小熊座图像）
\end{itemize}